%! Periodic Phenomena — Unit 6, Lesson 6.1 — Group Activity
\documentclass[10pt]{article}
\usepackage{precalculus-article}
\usepackage{precalculus-boxes}
\newcommand{\TallMath}[1]{$\displaystyle #1\rule[-1.4em]{0pt}{3.2em}$}

\begin{document}

\pageheader{Unit 6, Lesson 6.1}{Group Activity}
\namepartnerperiod

\vspace{0.10in}

%----------------------------------------------------------------------
% Context
%----------------------------------------------------------------------
\begin{scenariobox}[Context: Periodic Patterns in the World]{plumlight}
Periodic relationships appear throughout nature, engineering, and everyday
life. In this activity your group will investigate several real-world
contexts, determine whether each is periodic, estimate periods, and describe
key characteristics. Choose \textbf{one tier} below based on your readiness.
Work together --- everyone should be ready to explain your group's reasoning.
\end{scenariobox}

\vspace{0.08in}

%----------------------------------------------------------------------
% Tier R
%----------------------------------------------------------------------
\begin{tcolorbox}[colback=white, colframe=black!40,
    title=\textbf{Tier R --- Reinforce},
    fonttitle=\bfseries, arc=2mm, left=3mm, right=3mm, top=2mm, bottom=2mm]

\textbf{Context:} The table shows the average daily high temperature
(${}^\circ$F) in Portland, OR by month over two years.

\vspace{4pt}
\begin{center}
\renewcommand{\arraystretch}{1.35}
\small
\begin{tabular}{|c|c|c|c|c|c|c|c|c|c|c|c|c|}
\hline
\textbf{Month} & 1 & 2 & 3 & 4 & 5 & 6 & 7 & 8 & 9 & 10 & 11 & 12 \\
\hline
\textbf{Year 1} & 47 & 52 & 58 & 63 & 70 & 76 & 82 & 82 & 74 & 63 & 52 & 47 \\
\hline
\textbf{Month} & 13 & 14 & 15 & 16 & 17 & 18 & 19 & 20 & 21 & 22 & 23 & 24 \\
\hline
\textbf{Year 2} & 47 & 52 & 58 & 63 & 70 & 76 & 82 & 82 & 74 & 63 & 52 & 47 \\
\hline
\end{tabular}
\end{center}

\vspace{6pt}
\textbf{1.}\quad Is this relationship periodic? Circle one: \quad \textbf{YES} \quad / \quad \textbf{NO}

\vspace{4pt}
Explain how you know using the output values in the table:
\writelines{2}

\vspace{4pt}
\textbf{2.}\quad Estimate the period. The output pattern restarts at month
\underline{\hspace{1.2cm}}, so the period is $k =$ \underline{\hspace{1.2cm}} months.

\vspace{4pt}
\textbf{3.}\quad Describe what happens in \emph{one} cycle (months 1--12).
Use complete sentences.

\vspace{4pt}
\begin{itemize}[leftmargin=1.4em, itemsep=4pt]
  \item The temperature is \emph{increasing} from month
    \underline{\hspace{0.8cm}} to month \underline{\hspace{0.8cm}}.
  \item The temperature is \emph{decreasing} from month
    \underline{\hspace{0.8cm}} to month \underline{\hspace{0.8cm}}.
  \item The maximum temperature in one cycle is approximately
    \underline{\hspace{1.5cm}}${}^\circ$F, occurring at month
    \underline{\hspace{0.8cm}}.
  \item The minimum temperature in one cycle is approximately
    \underline{\hspace{1.5cm}}${}^\circ$F, occurring at month
    \underline{\hspace{0.8cm}}.
\end{itemize}

\vspace{4pt}
\textbf{4.}\quad Will the same maximum and minimum values appear in Year 3
(months 25--36)? Explain why or why not.
\writelines{2}

\end{tcolorbox}

\vspace{0.08in}

%----------------------------------------------------------------------
% Tier A
%----------------------------------------------------------------------
\begin{tcolorbox}[colback=white, colframe=black!40,
    title=\textbf{Tier A --- Apply},
    fonttitle=\bfseries, arc=2mm, left=3mm, right=3mm, top=2mm, bottom=2mm]

\textbf{Context:} A water wheel has a radius of 10 feet and rotates at a
constant rate, completing one full revolution every 8 seconds. A red paint
mark on the rim starts at the 3 o'clock position (height $= 0$ ft relative
to the axle). As it rotates counterclockwise, the height changes as follows
over one full revolution:

\vspace{4pt}
\begin{center}
\renewcommand{\arraystretch}{1.35}
\begin{tabular}{|c|c|c|c|c|c|}
\hline
Position in cycle & Start & $\tfrac{1}{4}$ turn & $\tfrac{1}{2}$ turn & $\tfrac{3}{4}$ turn & Full turn \\
\hline
Time $t$ (sec) & 0 & 2 & 4 & 6 & 8 \\
\hline
Height $h$ (ft) & 0 & 10 & 0 & $-10$ & 0 \\
\hline
\end{tabular}
\end{center}

\vspace{6pt}
\textbf{1.}\quad Build a table of height values for \textbf{two full periods}
($t = 0$ through $t = 16$, every 2 seconds).

\vspace{4pt}
\begin{center}
\renewcommand{\arraystretch}{1.5}
\begin{tabular}{|c|c|c|c|c|c|c|c|c|c|}
\hline
$t$ (sec) & 0 & 2 & 4 & 6 & 8 & 10 & 12 & 14 & 16 \\
\hline
$h(t)$ (ft) & & & & & & & & & \\
\hline
\end{tabular}
\end{center}

\vspace{6pt}
\textbf{2.}\quad Is this relationship periodic? Justify using your table.
\writelines{2}

\vspace{4pt}
\textbf{3.}\quad Estimate the period from your table. $k =$ \underline{\hspace{1.5cm}} seconds.

\vspace{4pt}
\textbf{4.}\quad Over the interval $0 \leq t \leq 16$, describe the intervals
on which $h$ is \emph{increasing} and \emph{decreasing}.

\vspace{4pt}
\begin{itemize}[leftmargin=1.4em, itemsep=3pt]
  \item Increasing: \writeline
  \item Decreasing: \writeline
\end{itemize}

\vspace{4pt}
\textbf{5.}\quad Compute the average rate of change of $h$ from $t=0$ to $t=2$.
Show your work.
\writelines{2}

\end{tcolorbox}

\vspace{0.08in}

%----------------------------------------------------------------------
% Tier E
%----------------------------------------------------------------------
\begin{tcolorbox}[colback=white, colframe=black!40,
    title=\textbf{Tier E --- Extend},
    fonttitle=\bfseries, arc=2mm, left=3mm, right=3mm, top=2mm, bottom=2mm]

\textbf{Context:} Two ocean buoys record wave heights (in meters) above a
baseline. Data are recorded every hour.

\vspace{4pt}
\textbf{Buoy A:}
\begin{center}
\renewcommand{\arraystretch}{1.3}
\small
\begin{tabular}{|c|c|c|c|c|c|c|c|c|c|c|c|c|c|}
\hline
$t$ (hr) & 0 & 1 & 2 & 3 & 4 & 5 & 6 & 7 & 8 & 9 & 10 & 11 & 12 \\
\hline
$A(t)$ (m) & 0 & 1.5 & 2 & 1.5 & 0 & $-1.5$ & $-2$ & $-1.5$ & 0 & 1.5 & 2 & 1.5 & 0 \\
\hline
\end{tabular}
\end{center}

\vspace{4pt}
\textbf{Buoy B:}
\begin{center}
\renewcommand{\arraystretch}{1.3}
\small
\begin{tabular}{|c|c|c|c|c|c|c|c|c|c|c|c|c|c|}
\hline
$t$ (hr) & 0 & 1 & 2 & 3 & 4 & 5 & 6 & 7 & 8 & 9 & 10 & 11 & 12 \\
\hline
$B(t)$ (m) & 0 & 3 & 4 & 3 & 0 & $-3$ & $-4$ & $-3$ & 0 & 3 & 4 & 3 & 0 \\
\hline
\end{tabular}
\end{center}

\vspace{6pt}
\textbf{1.}\quad Is $A$ periodic? Is $B$ periodic? For each, identify the
period or explain why it is not periodic.
\writelines{3}

\vspace{4pt}
\textbf{2.}\quad Compute the average rate of change of $A$ from $t=0$ to $t=2$.
\writelines{2}

\vspace{4pt}
\textbf{3.}\quad Compute the average rate of change of $B$ from $t=0$ to $t=2$.
\writelines{2}

\vspace{4pt}
\textbf{4.}\quad Which buoy's wave height is changing more rapidly during the
interval $0 \leq t \leq 2$? What does this mean physically?
\writelines{2}

\vspace{4pt}
\textbf{5.}\quad Using EK 3.1.B.3, explain why the rates of change you computed
in \#2 and \#3 will appear at corresponding positions in \emph{every} period
of $A$ and $B$.
\writelines{2}

\end{tcolorbox}

\end{document}
